%! Sine and Cosine Function Values — Unit 6, Lesson 6.4 — Warm-Up (Answer Key)
\documentclass[10pt]{article}
\usepackage{precalculus-article}
\usepackage{precalculus-key}   % pulls in -boxes; do NOT also load -boxes
\newcommand{\TallMath}[1]{$\displaystyle #1\rule[-1.4em]{0pt}{3.2em}$}

\begin{document}

\pageheader{Unit 6, Lesson 6.4}{Warm-Up --- Answer Key}
\namedateperiod

\vspace{0.12in}

\begin{notesbox}{Special Right Triangles \& Unit Circle Review}

\textbf{1.} An isosceles right triangle has hypotenuse 1. Find the length of each leg.
Show your work using the Pythagorean theorem.

\medskip
$a^2 + a^2 = 1^2 \Rightarrow 2a^2 = 1 \Rightarrow a = \dfrac{1}{\sqrt{2}} = \dfrac{\sqrt{2}}{2}$

\vspace{0.1in}

Each leg $= $ \ans{$\dfrac{\sqrt{2}}{2}$}

\vspace{0.18in}

\textbf{2.} A 30-60-90 triangle has hypotenuse 1. Use the properties of an equilateral
triangle to find the short leg and the long leg.

\medskip
Bisecting an equilateral triangle of side 1 gives short leg $= \tfrac{1}{2}$ and
long leg $= \sqrt{1^2 - (1/2)^2} = \sqrt{3/4} = \tfrac{\sqrt{3}}{2}$.

\vspace{0.1in}

Short leg $= $ \ans{$\dfrac{1}{2}$} \qquad Long leg $= $ \ans{$\dfrac{\sqrt{3}}{2}$}

\vspace{0.18in}

\textbf{3.} A point $P$ lies on the unit circle at angle $\theta = \dfrac{\pi}{2}$.
From the unit-circle definition, state the coordinates of $P$.

\medskip
At $\theta = \pi/2$, the terminal ray points straight up: $(\cos(\pi/2), \sin(\pi/2)) = (0, 1)$.

\vspace{0.1in}

$P = \Bigl($ \ans{$0$} $,$ \ans{$1$} $\Bigr)$

\vspace{0.18in}

\textbf{4.} Suppose $\cos\theta = \dfrac{1}{2}$ and $\theta$ is in Quadrant~I.
Use $x^2 + y^2 = 1$ to find $\sin\theta$.

\medskip
$\left(\tfrac{1}{2}\right)^2 + y^2 = 1 \Rightarrow y^2 = \tfrac{3}{4} \Rightarrow y = \tfrac{\sqrt{3}}{2}$
(positive because QI).

\vspace{0.1in}

$\sin\theta = $ \ans{$\dfrac{\sqrt{3}}{2}$}

\end{notesbox}

\end{document}
